\documentclass[11pt,a4paper]{article}

\usepackage[utf8]{inputenc}
\usepackage[T1]{fontenc}
\usepackage[margin=2.2cm]{geometry}
\usepackage{parskip}
\usepackage{titlesec}
\usepackage{enumitem}
\usepackage{booktabs}
\usepackage{xcolor}
\usepackage{tcolorbox}
\usepackage{hyperref}
\usepackage{fancyhdr}
\usepackage{listings}
\usepackage{longtable}
\usepackage{array}

\tcbuselibrary{skins,breakable}

\definecolor{accent}{RGB}{30,80,150}
\definecolor{softgray}{RGB}{245,246,248}
\definecolor{codebg}{RGB}{248,248,250}
\definecolor{warn}{RGB}{160,60,40}
\definecolor{good}{RGB}{30,120,70}

\titleformat{\section}{\Large\bfseries\color{accent}}{\thesection}{0.8em}{}
\titleformat{\subsection}{\large\bfseries}{\thesubsection}{0.6em}{}
\titleformat{\subsubsection}{\normalsize\bfseries\itshape}{\thesubsubsection}{0.5em}{}

\hypersetup{
  colorlinks=true, linkcolor=accent, urlcolor=accent,
  pdftitle={AgentIA --- Reflective Skill Optimisation: Proposal, Lineage and Schema},
  pdfauthor={AgentIA Engineering}
}

\pagestyle{fancy}
\fancyhf{}
\fancyhead[L]{\small\color{accent}AgentIA --- Reflective Skill Optimisation}
\fancyhead[R]{\small\color{accent}\thepage}
\renewcommand{\headrulewidth}{0.4pt}

\lstset{
  basicstyle=\ttfamily\small, backgroundcolor=\color{codebg},
  frame=single, rulecolor=\color{softgray}, breaklines=true,
  columns=fullflexible, keepspaces=true, showstringspaces=false,
  xleftmargin=0.5em, xrightmargin=0.5em,
}

\newtcolorbox{notebox}[1][]{colback=softgray,colframe=accent,boxrule=0.4pt,
  arc=2pt,left=6pt,right=6pt,top=4pt,bottom=4pt,breakable,#1}

\title{\vspace{-1.4cm}\textbf{Reflective Skill Optimisation}\\[0.25cm]
\large Proposal, research lineage, and implementation schema for a no-RL
self-improvement loop over the diagnostic channel}
\author{AgentIA Engineering}
\date{2026-09-29}

\begin{document}
\maketitle
\thispagestyle{fancy}

\begin{notebox}
\textbf{What this document is.} A proposal for making the AgentIA coding agent
more reliable using \textbf{reflective, no-RL optimisation}: a language model
reflects over recorded evidence and proposes edits to skill documents, which are
admitted only on measured contribution and evicted when they fall below an
evidence floor. \textbf{No weight updates, no training, no GPUs} --- the platform
has none, and the evidence says a reflective curator captures a measurable
fraction of what a trained one achieves.

It is written \emph{after} the first real run
(\texttt{agentia\_first\_baseline.tex}) and after the instrument shipped, so it
separates what has been measured from what is merely proposed. Where a claim is
not yet measured, it is labelled as such.
\end{notebox}

\tableofcontents
\vspace{0.5cm}

\section{The proposal in one page}

\begin{center}
\begin{tabular}{@{}p{0.24\textwidth}p{0.70\textwidth}@{}}
\toprule
\textbf{Element} & \textbf{Proposal} \\
\midrule
Objective & \textbf{Conformance}, measured by the deterministic rule instrument ---
not build success, which is a binary whole-corpus outcome that no skill edit can
move detectably at this sample size. \\
Proposal mechanism & A reflective LLM proposes a bounded edit set against skill
documents. No gradients, no fine-tuning, no RL. \\
Admission & Measured contribution, by leave-one-out: compare outcomes with the
skill present against outcomes with it absent. \\
Eviction & Remove a skill only when its measured contribution is \textbf{below a
stated floor}. ``No proven benefit'' is not grounds for removal. \\
Guardrail & The existing strict-improvement gate is demoted from acceptance
criterion to regression veto. It keeps its job; it loses its vote. \\
Bound & A round considers a bounded number of skills, because attribution cost is
(multiplicative) skills $\times$ tasks. \\
Gate to build & The corpus report must first show the conformance measure
\emph{varies} across real sessions. If every session scores identically, there is
nothing to move. \\
\bottomrule
\end{tabular}
\end{center}

The loop, in the shape it would take:

\begin{lstlisting}
  session runs
      |
      v
  compliance gate evaluates each stage's candidate
      |                       \
      | (pass)                 \ (fail) -> rule_id, severity, artifact, stage
      v                         v
  artifacts persisted      DIAGNOSTIC CHANNEL          <-- exists (feature 015)
                                |
                                v
                     aggregate by rule across sessions
                     (which rules recur, where, how often)
                                |
                                v
                     reflective LLM proposes skill edits
                     targeting the RECURRING rules
                                |
                                v
                     candidate skill -> fresh sessions -> measure
                                |
                                v
                     admit only on measured contribution;
                     evict skills below the evidence floor
\end{lstlisting}

\section{Why the diagnostic channel precedes the loop}
\label{sec:channel}

This ordering is the proposal's load-bearing decision, and it is not a preference.
It follows from the single largest effect in the evidence base.

CoEvoSkills (\texttt{arXiv:2604.01687}) reports \textbf{+40.5 points} on
SkillsBench with a co-evolving diagnostic verifier. Its own ablation:

\begin{center}
\begin{tabular}{@{}lr@{}}
\toprule
\textbf{CoEvoSkills variant} & \textbf{Score} \\
\midrule
Full (co-evolving diagnostic verifier) & \textbf{71.1} \\
Opaque pass/fail oracle only & 41.1 \quad ($-30.0$) \\
No evolution at all & 42.4 \quad ($-28.7$) \\
\bottomrule
\end{tabular}
\end{center}

\textbf{A loop fed only a binary verdict performs no better than no loop at all.}
41.1 with a pass/fail oracle against 42.4 with nothing: the loop is not merely
weaker without structure, it is \emph{indistinguishable from absent}. The
diagnostic channel is therefore worth more than the loop, and building the loop
first would have produced a component that cannot be shown to do anything.

Three further results point the same way:

\begin{itemize}[leftmargin=1.4em]
  \item An LLM-as-judge is the instrument the literature most distrusts, and it
  fails concretely: in the direct SkillOpt-on-code evaluation
  (\texttt{arXiv:2609.12742}), a candidate that wrote \textbf{an empty file} scored
  \emph{higher} on the judge than the baseline. A scoring channel that rewards
  emptiness cannot steer anything.
  \item Hand-written rule quality is largely content-independent --- random rules
  and curated rules scored \textbf{63.8\%} each in the guardrails study
  (\texttt{arXiv:2604.11088v2}, 679 files, 25{,}532 rules) --- which is an argument
  for a \emph{deterministic} instrument rather than for cleverer prose.
  \item Rules about architecture measured \textbf{worst} of the rule types
  (53.4\% against 63.8\% for tool/process rules). Since AgentIA's whole subject is
  layered architecture, the signal must come from the platform's own gate, not from
  a rule recipe imported from elsewhere.
\end{itemize}

This is why feature 015 shipped the instrument and \emph{deferred} the optimiser.
The decision is recorded in \texttt{specs/015-diagnostic-skill-evolution/spec.md}
under ``Deferred: the optimizer''.

\section{The three conditions the evidence requires}
\label{sec:conditions}

\subsection{Condition 1 --- Per-skill contribution, by leave-one-out}

A whole-corpus pass-rate comparison cannot attribute an effect to a rule: with one
document, co-applied edits confound one another, and ``the output got better'' is
not ``this rule was responsible''. The primitive that resolves this is
training-free and computes

\begin{center}
\texttt{delta(task, skill) = r(task, skill) - r(task, absent)}
\end{center}

by replaying with the skill removed. This is decision \textbf{D3} in
\texttt{research.md}. Its cost is multiplicative --- skills $\times$ tasks --- which
is why the per-round skill bound (D10) is not a convenience but the thing that
keeps measurement affordable.

\subsection{Condition 2 --- Eviction with an evidence floor}

A skill library that can only grow degrades until injecting a skill is worse than
injecting nothing (\texttt{arXiv:2605.22148}). The repair is contribution-based
eviction, and the floor is what stops measurement noise from driving removals:
``we could not prove it helps'' and ``it hurts'' are different findings.

The direction is corroborated by operation ablations
(\texttt{arXiv:2605.29440}):

\begin{center}
\begin{tabular}{@{}lr@{}}
\toprule
\textbf{Operation set} & \textbf{Score} \\
\midrule
Add only & 47.0 \\
Add $+$ remove & 48.3 \\
Add $+$ rewrite & 53.5 \\
All three & \textbf{59.0} \\
\bottomrule
\end{tabular}
\end{center}

Add-only editing --- feature 014's behaviour --- is the worst configuration. This is
decision \textbf{D4}.

\subsection{Condition 3 --- Diagnostics richer than pass/fail}

The channel must carry \emph{structure}: which rule, on which artifact, in which
stage, at which severity. A score says a session went badly; it does not say what
to change. \texttt{rule\_histogram} is the part a curator can act on, and
per-stage attribution converts a session total into a profile. This is decision
\textbf{D7}; the implementation reads the verdicts the platform \emph{already}
computes rather than re-evaluating anything.

\begin{notebox}[colframe=warn]
\textbf{The honest limit of this condition.} Per-stage attribution adds
\emph{resolution}, not \emph{evidence}. Re-checking the same task repeatedly raises
the session count but not the number of distinct tasks, and the channel is a fixed
hand-written rule set: a defect class nobody encoded reads clean. It measures
\textbf{conformance, not correctness} --- a well-layered but broken service scores
full marks.
\end{notebox}

\section{Research lineage}
\label{sec:lineage}

Three lineages, in ascending order of how directly they bear on this proposal.

\subsection{Reflective text-space optimisation (no weights)}

GEPA (\texttt{arXiv:2507.19457}) reports \textbf{+6\% average, up to +20\%} over
GRPO with \textbf{up to 35$\times$ fewer rollouts}. SkillOpt
(\texttt{arXiv:2605.23904}) reports best-or-tied in \textbf{52 of 52} cells,
\textbf{+5.4} over its strongest per-cell competitor (82.3 vs 76.9), across seven
models at roughly \textbf{+17.6} average. Both work on text, not weights.

\subsection{A trained curator, and the fraction reflection recovers}

CODESKILL (\texttt{arXiv:2605.25430}) is the closest controlled comparison, and it
is the source of the number this proposal leans on:

\begin{center}
\begin{tabular}{@{}lr@{}}
\toprule
\textbf{Curator} & \textbf{Score} \\
\midrule
No skill management & 29.32 \\
Prompt-based skill management (GPT-5.4-mini) & 35.25 \quad ($+5.93$, 54\% of headline) \\
SFT curator & 34.54 \quad ($+5.22$, 47\% of headline) \\
CODESKILL (RL curator) & \textbf{40.35} \quad ($+11.03$ headline) \\
\bottomrule
\end{tabular}
\end{center}

\textbf{A prompt-based curator with no training, no iteration and no verifier
feedback captures roughly half the gain of the trained RL curator.} That is the
measured floor reflection starts from --- not zero. The RL half ($+5.1$ to $+5.8$)
is the part this platform cannot reach, and it is honest to say so: the trained run
cost approximately \textbf{920 H100-GPU-hours}.

\subsection{Null results, which are part of the lineage}

\begin{itemize}[leftmargin=1.4em]
  \item The direct test of SkillOpt on repository-level code
  (\texttt{arXiv:2609.12742}) reports GEPA docs \textbf{+4.9\,pp} and SkillOpt
  \textbf{+0.1\,pp}, \emph{neither significant}, over 20--26 held-out tasks per repo.
  \item Gloaguen et al.\ (\texttt{arXiv:2602.11988}) found \textbf{no general gain}
  and \textbf{+20\% cost}; the v2 figure is ``7\%'' with a ``does not generally
  improve'' verdict.
  \item LLM-written skills measured \textbf{+0.0\,pp} against \textbf{+16.2\,pp} for
  human-written ones in the Ratchet study.
\end{itemize}

The reports themselves disagree on how to frame these; the fact sheet records that
report~1 calls the repo-level question untested while report~2 reports the direct
null. For this proposal the reading is simple: \textbf{the effect is real in
text-space but bounded, and it is procedural --- consistency, not capability.}

\subsection{Where the levers rank}

The two reports rank the available levers, and they differ in a way this proposal
must respect:

\begin{center}
\begin{tabular}{@{}p{0.30\textwidth}p{0.30\textwidth}p{0.32\textwidth}@{}}
\toprule
& \textbf{Report 1 (capability)} & \textbf{Report 2 (zero-compute EV)} \\
\midrule
Best & model capability & verification-guided best-of-$k$ \\
then & harness / interface design & executable house standard \\
then & verification-guided selection & conformance score as metric \\
then & RL post-training & harness / pipeline \\
then & RL-trained skill banks & retrieval \\
then & \textbf{optimised text skills} & bounded self-repair \\
worst & hand-written rules & \textbf{skill-document optimisation} \\
\bottomrule
\end{tabular}
\end{center}

\noindent On capability, text skills rank above hand-written rules. On expected
value per unit of compute, skill-document optimisation ranks \textbf{last}. This
proposal is therefore explicitly the \emph{last} lever, and it says so rather than
implying otherwise. It is being built because the mechanism is cheap
\emph{once the evidence exists} --- and because the evidence is the thing the
platform just learned it did not have.

\section{Implementation schema}
\label{sec:schema}

\subsection{The channel, end to end}

\begin{lstlisting}
  stage candidate
        |
        v
  compliance gate (security_service + test_analysis_service, merged)
        |  LOCAL violation  -> reject, consume a correction attempt
        |  ACCUMULATED      -> record, never reject
        v
  accepted artifact set --> persisted to the workspace
        |
        v
  session terminal ------> diagnose(artifacts)                 [instrument]
        |                       |
        |                       +--> score, raw_penalty, density
        |                       +--> rule_histogram, findings[]
        |                       +--> evaluable / not evaluable
        |                       +--> unverified (synthetic build)
        v
  stage_attribution(journal)
        |  reads initial_verdict per stage  (introduction, not persistence)
        |  per stage: outcome, requests, rule_histogram, passed|None
        v
  session_diagnostic_records  (one row per session)
        |
        v
  corpus_report: population, exclusions+reasons, conformance distribution,
                 correction effort, distinct tasks, cost if data exists
        |
        v
  [DEFERRED] skill_contribution: leave-one-out delta per skill
        |
        v
  [DEFERRED] reflect -> candidate skill -> fresh sessions -> admit/evict
\end{lstlisting}

\subsection{What exists, and what is deferred}

\begin{center}
\begin{tabular}{@{}p{0.33\textwidth}p{0.33\textwidth}p{0.26\textwidth}@{}}
\toprule
\textbf{Component} & \textbf{Where} & \textbf{State} \\
\midrule
Deterministic instrument & \texttt{services/conformance\_diagnostics.py} & \textbf{built} \\
Per-stage attribution & same, reading the stage journal & \textbf{built} \\
Session record (durable) & \texttt{models/diagnostics.py} & \textbf{built} \\
Corpus report & \texttt{services/corpus\_report.py} & \textbf{built} \\
Batch driver & \texttt{scripts/run\_corpus\_baseline.py} & \textbf{built} \\
Skill generator (reflection) & \texttt{scripts/skillopt/reflect.py} & exists, \textbf{fed the wrong signal} \\
\textbf{Per-skill contribution} & \texttt{services/skill\_contribution.py} & \textbf{deferred} \\
\textbf{Eviction with a floor} & \texttt{scripts/skillopt/currency.py} & \textbf{deferred} \\
\bottomrule
\end{tabular}
\end{center}

The reflector currently receives a derived build exit code, a terminal status and
artifact paths --- the ``opaque verdict'' of \S\ref{sec:channel}. It should receive
the rule histogram: which rules fired, how often, at which stage. The channel
exists; the wiring does not. That is a small change to a built component, and it is
deferred only because the channel's variance across real sessions is still
unmeasured.

\section{What is measured today}
\label{sec:measured}

\subsection{The first real run}

From \texttt{agentia\_first\_baseline.tex}: one session, blueprint \texttt{minimal},
provider \texttt{deepseek}, model \texttt{deepseek-flash}.

\begin{center}
\begin{tabular}{@{}lr@{}}
\toprule
\textbf{Measure} & \textbf{Value} \\
\midrule
Model calls & 15 \\
Input tokens & 26{,}351 \\
Output tokens & 139{,}679 \\
Cost & \$0.085295889 \quad (15 of 15 calls priced) \\
Terminal state & BLOCKED \\
\bottomrule
\end{tabular}
\end{center}

Every stage consumed its entire correction budget. The first four were accepted on
their third attempt; TEST never converged. This is \emph{systematic procedural
non-compliance} --- exactly the failure mode a reflective loop targets, and
precisely the signal the compliance gate already computes.

\subsection{The artifact question, re-measured}

The baseline document recorded a diagnostic record with
\texttt{artifact\_count = 0}, \texttt{evaluable = false} on a session whose first
four stages had persisted artifacts, and named that the critical path. That has now
been traced to its end, and the conclusion is different from the one the record
implied.

\begin{center}
\begin{tabular}{@{}p{0.42\textwidth}p{0.50\textwidth}@{}}
\toprule
\textbf{Question} & \textbf{Answer, and how it was established} \\
\midrule
Are artifacts lost on the stage-runner exhaustion path? & \textbf{No.} A scripted
model reproducing the baseline's exact pattern (two rejections then an acceptance
for each of four stages, three rejections for TEST) returns all seven persisted
artifacts in the final state --- through the stage loop, through the compiled graph,
and through the production \texttt{stream}+update accumulator. \\
Why did the record read zero? & It was \textbf{fabricated test output}. The single
stored record held five EXHAUSTED stages with empty rule histograms --- the shape
the test fixtures wrote under the production session key, which is the collision
commit \texttt{dbcfa2e} fixed. \\
Is a real blocked session diagnosable now? & \textbf{Yes.} The reproduction records
\texttt{artifact\_count = 7}, \texttt{evaluable = true}, with the four corrected
stages carrying the rules that rejected their first candidate. \\
\bottomrule
\end{tabular}
\end{center}

A regression test now pins each of those three claims
(\texttt{test\_diagnostic\_record.py}), so the ``not yet reproduced'' note in
\texttt{dbcfa2e} is closed.

One genuine defect \emph{was} found by tracing this, and it matters more than the
one that was assumed. \texttt{stage\_attribution} read
\texttt{verdict.get("passed", not violations)} on an absent verdict, so a stage
whose candidate never parsed was reported \texttt{passed = true}. That rendered a
session in which \emph{every} stage was rejected as one in which every stage
passed --- the conflation FR-004 forbids, one level down, and the exact input an
optimiser would learn nothing from. A stage with no verdict now reports
\texttt{passed = null}: absence, not success.

\section{The gate on building the loop}
\label{sec:gate}

\begin{notebox}[colframe=warn]
\textbf{Do not build the optimiser until the corpus report shows the conformance
measure VARIES across real sessions.} If every recorded session scores identically,
there is nothing for a loop to move, and the effort belongs on a different lever.
\end{notebox}

Even once it varies, the arithmetic is unforgiving, and it is arithmetic rather
than pessimism. For a paired comparison, the exact floor at five tasks cannot reach
significance under \emph{any} outcome:

\begin{center}
\begin{tabular}{@{}lrl@{}}
\toprule
\textbf{Distinct tasks} & \textbf{Best attainable $p$} & \textbf{Verdict} \\
\midrule
4 & 0.125 & cannot reach $p < 0.05$ \\
5 & 0.0625 & cannot reach $p < 0.05$ \\
6 & 0.03125 & first size that can \\
\bottomrule
\end{tabular}
\end{center}

Against a single-run pass@1 variance measured at \textbf{2.2--6.0\,pp}, five
blueprints support no inference at all: one task moves the pass rate 25 points. The
power analysis is equally blunt: roughly \textbf{120--200 paired tasks} are needed
to detect a 10\,pp effect. The platform has five.

\subsection{So the decision is: grow the corpus, or do not build the loop}

In dependency order:

\begin{enumerate}[leftmargin=1.4em]
  \item \textbf{Grow the held-out corpus.} Six is the arithmetic minimum; the power
  target is two orders of magnitude beyond it. This is the highest-value change and
  it is cheap.
  \item \textbf{Run real sessions with the container runtime reachable}, so the
  diagnostic records accumulate and the variance question can actually be answered.
  \item \textbf{Feed the reflector the rule histogram} instead of the build exit
  code.
  \item \textbf{Build leave-one-out contribution}, then \textbf{eviction with a
  floor}.
  \item \textbf{Only then} run the loop, scored on conformance.
\end{enumerate}

\section{Honest limits}
\label{sec:limits}

\begin{center}
\begin{tabular}{@{}p{0.30\textwidth}p{0.64\textwidth}@{}}
\toprule
\textbf{Limit} & \textbf{Detail} \\
\midrule
Reflection's ceiling is unmeasured here & It captures $\approx$54\% of a trained
curator's gain \emph{on CODESKILL's benchmarks}. How much it recovers on repository-level
Spring Boot code is not stated anywhere in the reports. \\
The direct repo-level test is null & SkillOpt on repository-level code measured
\textbf{+0.1\,pp}, not significant. This proposal is betting on conformance and
cost, where the evidence is better, not on resolve rate. \\
The measure is conformance, not correctness & A well-layered but broken service
scores full marks. A defect class nobody encoded reads clean. \\
Stage-local findings are blocked before persistence & The findings that survive
into a saved artifact set come from whole-project rules only --- a narrow signal,
and its real-world variance is exactly what is still unmeasured. \\
Small $n$ & Five blueprints. No significance is attainable. Repeated runs of one
task are not independent observations and are never counted as a larger sample. \\
No trainable environment & No GPUs, no fine-tuning, no RL. The trained half of the
CODESKILL result is out of reach by construction, not by choice. \\
The corpus currently holds no real records & The store held one record, and it was
test output. A measurement store that cannot be told apart from its own fixtures is
worthless; that is fixed, but it means the variance question is still open. \\
\bottomrule
\end{tabular}
\end{center}

\section{Decision record}

\begin{center}
\begin{tabular}{@{}p{0.10\textwidth}p{0.44\textwidth}p{0.38\textwidth}@{}}
\toprule
\textbf{Ref} & \textbf{Decision} & \textbf{Reason} \\
\midrule
D1 & Optimise conformance, not build success & build success is binary and
unmeasurable at $n=5$ \\
D2 & The instrument is deterministic, not an LLM judge & a drifting judge makes a
measured improvement unattributable \\
D3 & Contribution by leave-one-out, per skill & the only training-free primitive
that converts ``it got better'' into ``this rule did it'' \\
D4 & Remove only below an evidence floor & ``not proven to help'' $\neq$ ``hurts'';
noise must not drive removals \\
D5 & Reflect; never train & no compute for weight updates; reflection has a
measured floor \\
D6 & Normalise the measure for artifact-set size & a raw count rewards generating
less code \\
D7 & Attribution from the existing stage journal & it is already computed and free
to keep \\
D8 & Insufficiency is a reportable outcome & a point estimate without its power is
noise reported as a result \\
D9 & Demote the gate to a guardrail & it is a good regression veto and a bad
acceptance criterion \\
D10 & Bound the skills per round & attribution cost is skills $\times$ tasks \\
\bottomrule
\end{tabular}
\end{center}

\begin{notebox}
\textbf{The one-sentence position.} Text-space skill optimisation is real, bounded,
and ranked last per unit of compute; its diagnostic channel is worth more than its
loop, and the platform now has the channel, the record, and the report --- so the
next honest step is to grow the corpus and measure whether there is anything to
optimise, not to build the optimiser.
\end{notebox}

\end{document}
